% =====================================================================
%  Formatting preamble for the Pandoc -> xelatex build.
%  Consumed via --include-in-header (see defaults.yaml); injected after
%  Pandoc's own template setup, so it loads on top of it.
%
%  Font stack adapted for xelatex. The original pdflatex stack
%  (libertinus + libertinust1math + the inconsolata package) is replaced by
%  libertinus-otf (Libertinus Serif text + Libertinus Math via unicode-math)
%  plus a fontspec monospace, because libertinust1math is incompatible with
%  xelatex (it breaks \boldsymbol). Document geometry, colored links, the
%  document class, and section numbering live in metadata.yaml.
% =====================================================================

% ---- fonts (xelatex) -------------------------------------------------
\usepackage{libertinus-otf}                              % Libertinus Serif + Libertinus Math
\setmonofont{Inconsolatazi4-Regular.otf}[Scale=0.96]     % Inconsolata monospace
\providecommand{\bm}[1]{\boldsymbol{#1}}                 % map \bm to unicode-math \boldsymbol

% ---- pandoc-crossref multi-image figures -----------------------------
% A `::: {#fig:x}` div holding several images emits a pandoccrossrefsubfigures
% environment that crossref (0.3.x) expects the document to define. We use no
% sub-captions, so a plain centered figure float is the correct wrapper.
\newenvironment{pandoccrossrefsubfigures}{\begin{figure}\centering}{\end{figure}}

% ---- tables ----------------------------------------------------------
\usepackage{etoolbox}
\AtBeginEnvironment{longtable}{\small}   % render the converted longtables a touch smaller
\usepackage{calc}                        % column-width arithmetic used by converted tables
\usepackage{xurl}                        % allow long URLs to break across lines

% ---- helper macros emitted by Pandoc fragments (providecommand: no clash) ----
\providecommand{\tightlist}{\setlength{\itemsep}{0pt}\setlength{\parskip}{0pt}}
\providecommand{\real}[1]{#1}
\providecommand{\passthrough}[1]{#1}

% ---- framed placeholder for figures to be supplied later (used in Appendix C) ----
\newcommand{\imgplaceholder}[1]{%
  \par\medskip
  \begin{center}%
  \fbox{\parbox[c][3cm][c]{0.82\linewidth}{\centering\itshape
  Figure placeholder: image \texttt{#1} to be supplied.}}%
  \end{center}%
  \par\medskip}

\emergencystretch=2em   % last-resort stretch to curb overfull prose lines
